\documentclass[a4paper]{article}

\usepackage[fontset=fandol]{ctex}
\usepackage{amsmath, amssymb}
\usepackage{graphicx}
\usepackage[margin=2.3cm]{geometry}
\usepackage[most]{tcolorbox}
\usepackage{etoolbox}
\usepackage{listings}
\usepackage{hyperref}
\usepackage{xurl}
\usepackage{booktabs}
\usepackage{longtable}
\usepackage{tabularx}
\usepackage{array}
\usepackage{float}
\usepackage{enumitem}
\usepackage{tikz}

\hypersetup{
    colorlinks=true,
    linkcolor=blue!50!black,
    urlcolor=blue!60!black,
    citecolor=blue!50!black
}

\newtcolorbox{findingbox}[1]{
    enhanced,
    colback=green!5!white,
    colframe=green!45!black,
    colbacktitle=green!45!black,
    coltitle=white,
    fonttitle=\bfseries,
    title=#1,
    sharp corners
}

\newtcolorbox{evidencebox}[1]{
    enhanced,
    colback=blue!4!white,
    colframe=blue!65!black,
    colbacktitle=blue!65!black,
    coltitle=white,
    fonttitle=\bfseries,
    title=#1,
    sharp corners
}

\newtcolorbox{riskbox}[1]{
    enhanced,
    colback=red!4!white,
    colframe=red!70!black,
    colbacktitle=red!70!black,
    coltitle=white,
    fonttitle=\bfseries,
    title=#1,
    sharp corners
}

\newtcolorbox{methodbox}[1]{
    enhanced,
    colback=yellow!9!white,
    colframe=yellow!60!black,
    colbacktitle=yellow!60!black,
    coltitle=black,
    fonttitle=\bfseries,
    title=#1,
    sharp corners
}

\lstset{
    language=bash,
    basicstyle=\ttfamily\small,
    keywordstyle=\color{blue},
    stringstyle=\color{red!60!black},
    commentstyle=\color{green!50!black},
    breaklines=true,
    frame=single,
    numbers=left,
    numberstyle=\tiny\color{gray},
    captionpos=b,
    extendedchars=false
}

\newcolumntype{Y}{>{\raggedright\arraybackslash}X}
\setlist[itemize]{leftmargin=1.8em}
\setlength{\parskip}{0.45em}

\newcommand{\reporttitle}{video-note 测试体系现状审计}
\newcommand{\reportsubtitle}{runtime repo 与 thin wrapper 的覆盖盘点、缺口分析与增量测试路线图}
\newcommand{\reportauthors}{Codex}
\newcommand{\reportdate}{2026-03-26}
\newcommand{\researchquestion}{我们现在到底有哪些测试，它们实际保护了什么风险，以及下一批最值得增加哪些测试？}
\newcommand{\reportscope}{以 \path{/home/fanghaotian/Desktop/src/video-note-pipeline} 为主，结合 \path{/home/fanghaotian/Desktop/src/agent-basic-skill/skills/video-note-render-pdf} 的 wrapper 边界。证据窗口截止 2026-03-26。}
\newcommand{\reportaudience}{video-note runtime / wrapper 的维护者，以及准备继续扩展 subtitle-first 与 live E2E 回归的人。}
\newcommand{\sourcewindow}{本地源代码、文档、pytest 实测日志与最新 live E2E 产物。}

\begin{document}

\begin{titlepage}
\centering
{\Large 深度研究报告\par}
\vspace{1.0cm}
{\huge\bfseries \reporttitle\par}
\vspace{0.5cm}
{\Large \reportsubtitle\par}
\vspace{0.9cm}
{\large \reportauthors\par}
\vspace{0.2cm}
{\large \reportdate\par}
\vfill
\begin{tcolorbox}[width=0.94\textwidth, colback=black!2!white, colframe=black!60, sharp corners]
\textbf{核心问题}：\researchquestion\par
\textbf{研究范围}：\reportscope\par
\textbf{目标读者}：\reportaudience\par
\textbf{证据窗口}：\sourcewindow\par
\end{tcolorbox}
\end{titlepage}

\tableofcontents
\newpage

\section{执行摘要}

\begin{findingbox}{结论 1：现有测试资产并不薄弱，问题在于分布不均}
截至本次审计，\path{/home/fanghaotian/Desktop/src/video-note-pipeline} 的运行时测试实测为 \textbf{43/43 通过}，\path{/home/fanghaotian/Desktop/src/agent-basic-skill} 中与 \texttt{video-note-render-pdf} 直接相关的 wrapper 测试为 \textbf{17/17 通过}，最新 live E2E canary 为 \textbf{3/3 case 通过}。因此当前状态不是“几乎没测试”，而是已经形成了 \textbf{单元/契约层、命令编排层、wrapper 解析层、live E2E 层} 四层结构。
\end{findingbox}

\begin{findingbox}{结论 2：最强保护层是命令编排与工件契约，最弱保护层是平台真实边界}
当前最扎实的部分，是 \texttt{prepare / probe / transcript / overview / build} 这条命令编排链条，以及 case artifact contract 的稳定性。当前最弱的部分，是 \textbf{adapter 真正面对平台 API 漂移的解析层}、\textbf{transcript parser 本身}、\textbf{CLI 错误语义}，以及 \textbf{真实 subtitle 路径的 live E2E}。这意味着系统内部 glue code 已经相对稳，但最靠近外部依赖和最终用户入口的地方仍然偏薄。
\end{findingbox}

\begin{findingbox}{结论 3：下一批最值得加的不是“更多同类 mock 测试”，而是补齐真实边界}
按照本报告的优先级模型，最值得新增的测试依次是：
\begin{itemize}
\item 真实平台字幕路径的 live canary：至少一条公开字幕 case 和一条 cookies 字幕 case
\item \texttt{transcripts.py} 的 direct parser tests
\item \texttt{adapters/*.py} 与 \texttt{probes/subtitles.py} 的 direct contract tests
\item \texttt{cli.py} 的命令行语义测试
\item wrapper 与 runtime 的 contract sync tests
\end{itemize}
\end{findingbox}

\begin{riskbox}{重要限定}
本次审计是在 \path{/home/fanghaotian/Desktop/src/video-note-pipeline} 的 \textbf{dirty worktree} 上完成的；因此结论反映的是当前本地工作副本，而不是某个已打 tag 的发布版本。相关状态见附件 \path{runtime-git-status.txt}。
\end{riskbox}

\section{研究范围与方法}

\begin{methodbox}{最小合理假设}
用户没有单独限定“只看 runtime”或“只看 wrapper”，因此本报告采取最小合理范围：以 \texttt{video-note-pipeline} 为主分析对象，把 \texttt{video-note-render-pdf} wrapper 视为跨仓边界，并明确区分“runtime 内部测试”和“wrapper 边界测试”。
\end{methodbox}

本次审计实际使用了三类证据：

\begin{itemize}
\item \textbf{源码与测试文件}：例如 \path{/home/fanghaotian/Desktop/src/video-note-pipeline/tests/test_commands.py}、\path{/home/fanghaotian/Desktop/src/video-note-pipeline/tests/test_cookies.py}、\path{/home/fanghaotian/Desktop/src/agent-basic-skill/tests/test_wrapper_resolution.py}
\item \textbf{文档与契约}：例如 \path{/home/fanghaotian/Desktop/src/video-note-pipeline/README.zh-CN.md}、\path{/home/fanghaotian/Desktop/src/agent-basic-skill/skills/video-note-render-pdf/references/adapter-contract.md}
\item \textbf{实测结果}：本次审计保存的 \path{runtime-pytest.log}、\path{wrapper-pytest.log}，以及已有的 \path{/home/fanghaotian/Desktop/src/video-note-pipeline/reports/live-e2e/runs/bilibili-core-base-model-rerun/summary.json}
\end{itemize}

\begin{lstlisting}[caption={本次审计实际执行的关键验证命令}]
uv run pytest /home/fanghaotian/Desktop/src/video-note-pipeline/tests
uvx pytest tests/test_wrapper_resolution.py tests/test_install_skill.py tests/test_skill_repo.py
\end{lstlisting}

为了给“应该先加哪些测试”排序，我使用了一个显式的优先级函数：

$$
P = 2I + B + D - C
$$

\begin{itemize}
\item $P$：建议测试的优先级分数，越高越该先做
\item $I$：如果该区域出问题，对主链路和归档质量的影响，取值 $1$ 到 $5$
\item $B$：是否跨越 repo、平台或外部工具边界，取值 $1$ 到 $3$
\item $D$：外部漂移风险，例如 yt-dlp、平台字幕接口、浏览器 cookies 行为，取值 $1$ 到 $3$
\item $C$：新增该测试的工程成本，取值 $1$ 到 $3$
\end{itemize}

这个公式不是现成标准，而是本次审计的 \textbf{方法学工具}。它的目的，是避免建议列表退化成“把所有未测试模块都补一遍”的平均主义。

\section{当前测试版图}

\subsection{四层测试栈}

\begin{figure}[H]
\centering
\begin{tikzpicture}[font=\small]
\node[draw, rounded corners, fill=green!8, minimum width=12.5cm, minimum height=1.35cm] (e2e) at (0,0) {\textbf{Layer 4: Live E2E canary} \quad 3/3 case 通过 \quad 当前仅覆盖 Bilibili + ASR fallback};
\node[draw, rounded corners, fill=blue!7, minimum width=12.5cm, minimum height=1.35cm] (runtime) at (0,-1.9) {\textbf{Layer 3: Runtime pytest} \quad 43/43 通过 \quad 以命令编排、artifact contract、cookies、mode routing 为主};
\node[draw, rounded corners, fill=yellow!10, minimum width=12.5cm, minimum height=1.35cm] (wrapper) at (0,-3.8) {\textbf{Layer 2: Wrapper pytest} \quad 17/17 通过 \quad 以路径解析、安装器、skill 包装完整性为主};
\node[draw, rounded corners, fill=gray!10, minimum width=12.5cm, minimum height=1.35cm] (docs) at (0,-5.7) {\textbf{Layer 1: 文档与契约} \quad SKILL / adapter-contract / case-bundle-contract / README};
\draw[->, thick] (docs.north) -- (wrapper.south);
\draw[->, thick] (wrapper.north) -- (runtime.south);
\draw[->, thick] (runtime.north) -- (e2e.south);
\end{tikzpicture}
\caption{当前测试栈示意图。}
\end{figure}

\noindent\small\emph{来源：\path{runtime-pytest.log}、\path{wrapper-pytest.log}、\path{/home/fanghaotian/Desktop/src/video-note-pipeline/reports/live-e2e/runs/bilibili-core-base-model-rerun/summary.json}，作者自绘。}\normalsize

\subsection{Runtime repo：43 个 pytest 用例}

\begin{table}[H]
\centering
\small
\begin{tabularx}{\textwidth}{>{\raggedright\arraybackslash}p{2.6cm}cYY}
\toprule
测试文件 & 数量 & 主要保护的风险 & 当前仍明显缺的部分 \\
\midrule
\texttt{test\_paths.py} & 6 & URL 平台识别、video id 抽取、workspace/config root 解析 & host 变体与奇怪 URL 形式还不够多 \\
\texttt{test\_contracts.py} & 3 & artifact 索引、inventory 结果、case contract 稳定性 & 没有跨 repo 的 schema sync 保护 \\
\texttt{test\_cookies.py} & 9 & keyring 选择、cookies 导出、平台过滤、去重、状态提示 & 真实浏览器 decrypt 和真实登录态仍未进入自动化 \\
\texttt{test\_transcriber\_probe.py} & 4 & transcriber candidate 排序与 skip 逻辑 & 真正的 CUDA / MLX sample 行为仍主要靠本机人工验证 \\
\texttt{test\_frames.py} & 4 & montage 输出与 \texttt{recommended\_mode} 路由启发式 & 真实视频分布仍只覆盖有限 archetype \\
\texttt{test\_build.py} & 2 & build happy path、外部 \texttt{note.tex} 支持 & latexmk 失败、preview 分支、缺失 PDF 分支没单测 \\
\texttt{test\_commands.py} & 11 & 命令编排、stage 更新、subtitle/cookies/ASR fallback、manifest 更新 & 仍以 fake adapter 为主，外部平台解析被 mock 掉了 \\
\texttt{test\_live\_e2e.py} & 4 & manifest 解析、case evaluation、smoke note 生成 & 只测 harness 自身，不测真实网络视频 \\
\bottomrule
\end{tabularx}
\caption{runtime repo 的当前 pytest 结构。}
\end{table}

\noindent\small\emph{计数来源：\path{runtime-test-counts.csv} 与 \path{runtime-pytest.log}。}\normalsize

\subsection{Wrapper repo：17 个 pytest 用例}

\begin{table}[H]
\centering
\small
\begin{tabularx}{\textwidth}{>{\raggedright\arraybackslash}p{2.8cm}cYY}
\toprule
测试文件 & 数量 & 主要保护的风险 & 当前仍明显缺的部分 \\
\midrule
\texttt{test\_wrapper\_resolution.py} & 7 & runtime repo 与 workspace 路径解析优先级 & 只验证路径，不验证运行后 contract 是否匹配 \\
\texttt{test\_install\_skill.py} & 5 & 安装器能识别外部 runtime repo，并正确复制 wrapper skill & 不验证安装后的真实 runbook 与 runtime 兼容性 \\
\texttt{test\_skill\_repo.py} & 5 & skill 包装目录包含预期文件、wrapper 结构完整 & 不验证文档字段与 runtime 实现是否语义同步 \\
\bottomrule
\end{tabularx}
\caption{wrapper 侧与 \texttt{video-note-render-pdf} 直接相关的当前 pytest 结构。}
\end{table}

\noindent\small\emph{计数来源：\path{wrapper-test-counts.csv} 与 \path{wrapper-pytest.log}。}\normalsize

\subsection{Live E2E：当前已经覆盖的真实 case}

当前 live E2E manifest 位于 \path{/home/fanghaotian/Desktop/src/video-note-pipeline/live_e2e/bilibili-core.json}，最新一次完整结果位于 \path{/home/fanghaotian/Desktop/src/video-note-pipeline/reports/live-e2e/runs/bilibili-core-base-model-rerun/summary.json}。

\begin{table}[H]
\centering
\small
\begin{tabularx}{\textwidth}{>{\raggedright\arraybackslash}p{2.4cm}YYc}
\toprule
Case & 预期模式与覆盖价值 & 当前观测到的 transcript / subtitle 路径 & 结果 \\
\midrule
\texttt{BV1GZw9zcEdJ} & \texttt{visual-light}，代表轻量口播/低视觉密度 & \texttt{subtitle\_source = none}，\texttt{transcript\_source = asr:faster-whisper:cpu} & PASS \\
\texttt{BV1cptBzFE5Q} & \texttt{static-outline}，代表单画布大纲课件 & \texttt{subtitle\_source = none}，\texttt{transcript\_source = asr:faster-whisper:cpu} & PASS \\
\texttt{BV1G5AAzHEH3} & \texttt{board-heavy}，代表高视觉密度屏录/白板 & \texttt{subtitle\_source = none}，\texttt{transcript\_source = asr:faster-whisper:cpu} & PASS \\
\bottomrule
\end{tabularx}
\caption{当前 live E2E canary 已覆盖的三类 archetype。}
\end{table}

\noindent\small\emph{来源：\path{/home/fanghaotian/Desktop/src/video-note-pipeline/reports/live-e2e/runs/bilibili-core-base-model-rerun/summary.md}。}\normalsize

\begin{riskbox}{这里的关键限制}
这套 live E2E 虽然很有价值，但它目前 \textbf{只证明了 ASR fallback 主链路}。它还没有自动证明：
\begin{itemize}
\item 公开平台字幕路径在真实视频上持续可用
\item cookies 字幕路径在真实登录态下持续可用
\item YouTube 与 Bilibili 的平台差异在 live 环境中同时稳定
\end{itemize}
\end{riskbox}

\section{这些测试目前真正保护了什么}

\subsection{命令编排和工件契约已经比较稳}

\begin{evidencebox}{证据链}
\path{/home/fanghaotian/Desktop/src/video-note-pipeline/tests/test_commands.py} 覆盖了 \texttt{prepare / probe / transcript / overview / build / cookies-export} 六类命令的主要 happy path 与关键 fallback；\path{/home/fanghaotian/Desktop/src/video-note-pipeline/tests/test_contracts.py} 则确保 artifact 索引和 inventory 不会无声漂移。两者结合，意味着 case 生命周期的“骨架”已经被有效固定住。
\end{evidencebox}

如果只从归档风险看，当前最不容易被无声破坏的部分，是：

\begin{itemize}
\item case 目录结构与 artifact 命名
\item stage 状态推进
\item subtitle-first 在命令层的编排顺序
\item build 产物落位
\item \texttt{recommended\_mode} 写回 \texttt{preflight.json} 与 manifest
\end{itemize}

\subsection{mode routing 已经有“单测 + live case”的双重保护}

\begin{evidencebox}{证据链}
\path{/home/fanghaotian/Desktop/src/video-note-pipeline/tests/test_frames.py} 提供启发式单测；\path{/home/fanghaotian/Desktop/src/video-note-pipeline/reports/live-e2e/runs/bilibili-core-base-model-rerun/summary.json} 证明真实视频下三类 archetype 都能得到当前预期模式。
\end{evidencebox}

这意味着 \texttt{talking-head / visual-light / static-outline / board-heavy} 这套路由体系，已经不再只是 README 里的设计词汇，而是有运行时实证支撑的行为层输出。

\subsection{wrapper 侧保护的是“接线正确”，不是“语义正确”}

\begin{evidencebox}{证据链}
\path{/home/fanghaotian/Desktop/src/agent-basic-skill/tests/test_wrapper_resolution.py}、\path{/home/fanghaotian/Desktop/src/agent-basic-skill/tests/test_install_skill.py}、\path{/home/fanghaotian/Desktop/src/agent-basic-skill/tests/test_skill_repo.py} 证明 wrapper 能找到 runtime repo、能按安装器要求落盘、也包含所需参考文件。
\end{evidencebox}

但这些测试目前并不证明以下语义层断言：

\begin{itemize}
\item wrapper 文档里声明的必要 artifact 是否始终和 runtime 输出一致
\item wrapper 的 subtitle-first 文字约束是否与 runtime 的真实 fallback 行为持续同步
\item wrapper runbook 中的命令示例是否始终和 runtime CLI 参数保持一致
\end{itemize}

\section{当前明显未被充分保护的风险}

\subsection{未被 direct test 覆盖的关键模块}

根据附件 \path{source-module-test-map.txt}，以下模块目前没有独立的 direct test file：

\begin{table}[H]
\centering
\small
\begin{tabularx}{\textwidth}{>{\raggedright\arraybackslash}p{3.2cm}YY}
\toprule
模块 & 为什么重要 & 现在的问题 \\
\midrule
\texttt{adapters/base.py}、\texttt{adapters/youtube.py}、\texttt{adapters/bilibili.py} & 它们直接面对 yt-dlp 与 YouTubeTranscriptApi，是最容易被外部平台漂移打穿的边界 & 当前大部分命令测试都 fake 掉 adapter，因此平台解析层本身缺少稳定保护 \\
\texttt{probes/subtitles.py} & 这里编码了匿名字幕、cookies 字幕、失败原因拼接、selected payload 落盘等逻辑 & 虽然命令层间接覆盖到一部分，但没有更细粒度的 direct assertions \\
\texttt{transcripts.py} & 统一 transcript 三件套的 parser 与 writer 都在这里，是事实层核心格式化模块 & 目前没有 \texttt{test\_transcripts.py} 一类的 direct parser tests \\
\texttt{cli.py} & 这是最终用户真正会调用的命令入口 & 当前没有 direct CLI parser / exit code / JSON output tests \\
\texttt{utils/io.py}、\texttt{utils/process.py} & 是若干命令的底层依赖 & 风险较低，但一旦出问题会影响多个命令 \\
\bottomrule
\end{tabularx}
\caption{当前没有 direct test file 的主要模块。}
\end{table}

\noindent\small\emph{来源：\path{source-module-test-map.txt}。}\normalsize

\subsection{真实 subtitle-first 只在 mock 层被证明，没有在 live 层被证明}

这是当前最重要的测试缺口。原因不在于文档没写，而在于：

\begin{itemize}
\item wrapper 文档 \path{/home/fanghaotian/Desktop/src/agent-basic-skill/skills/video-note-render-pdf/SKILL.md} 明确要求先匿名字幕、再 cookies、最后 ASR
\item runtime 命令层的测试也已经验证了这种编排
\item 但 live E2E manifest 里的 3 个 case 全部走的是 \texttt{subtitle\_source = none} 与 \texttt{transcript\_source = asr:faster-whisper:cpu}
\end{itemize}

因此，当前最强的结论是：“ASR fallback 主链路被 live 证明了”；而不是“subtitle-first 的三层路径都已经被 live 证明了”。

\subsection{当前没有 CI 痕迹，意味着通过状态不是持续保证}

本次审计没有在两个仓库里发现 \path{.github/workflows/} 下的工作流文件，也没有发现 \path{Makefile}、\path{justfile}、\path{tox.ini} 或 \path{noxfile.py} 这类标准化测试入口。换句话说，当前测试是“可运行的”，但还不算“持续执行的”。这会直接降低新增测试的实际价值，因为没有自动门禁时，再好的测试也容易只在本机偶尔跑一次。

\section{应该增加什么测试}

\subsection{优先级排序}

\begin{table}[H]
\centering
\small
\begin{tabularx}{\textwidth}{c>{\raggedright\arraybackslash}p{4.2cm}cYY}
\toprule
优先级 & 建议测试 & 分数 $P$ & 推荐落点 & 为什么现在加 \\
\midrule
P0 & 真实公开字幕 canary + 真实 cookies canary & 11 & \texttt{live\_e2e/} manifest、\texttt{scripts/run\_live\_e2e.py} & 当前 live 只覆盖 ASR fallback，最重要的规范路径还没有 real-world 证明 \\
P0 & \texttt{transcripts.py} direct parser tests & 10 & 新增 \texttt{tests/test\_transcripts.py} & transcript 是事实层核心，但当前完全没有 direct tests \\
P0 & adapter contract tests & 10 & 新增 \texttt{tests/test\_adapters.py} & 这里最靠近 yt-dlp / YouTubeTranscriptApi，是平台漂移首个冲击面 \\
P1 & \texttt{probes/subtitles.py} direct tests & 9 & 新增 \texttt{tests/test\_subtitle\_probe.py} & 命令层已覆盖编排，但 probe 模块本体仍缺细粒度保护 \\
P1 & CLI 语义测试 & 8 & 新增 \texttt{tests/test\_cli.py} & 当前最终用户入口没有 direct tests，参数解析和 exit code 漂移无门禁 \\
P1 & wrapper/runtime contract sync tests & 8 & \texttt{agent-basic-skill/tests/} 新增 contract sync file & 现有 wrapper 测试只保结构，不保语义同步 \\
P1 & build failure / preview 分支测试 & 7 & 扩展 \texttt{tests/test\_build.py} & 当前 build 主要保护 happy path，失败路径仍偏薄 \\
P2 & resume / rebuild / broken artifact recovery tests & 7 & 新增 \texttt{tests/test\_commands\_recovery.py} 或扩展现有命令测试 & 这类问题在真实工作流里常见，但当前覆盖不足 \\
\bottomrule
\end{tabularx}
\caption{基于 $P = 2I + B + D - C$ 的新增测试排序。}
\end{table}

\subsection{建议 1：把 subtitle-first 真正跑成 live 回归}

这不是再加一个 mock 测试，而是把当前 live 套件从“只验证 ASR fallback”升级成“验证整个 subtitle-first 梯度”。具体建议如下：

\begin{itemize}
\item 增加一个 \textbf{公开视频字幕} case：例如稳定可访问的 YouTube 公共字幕视频，期望 \texttt{selected\_source = official-anonymous}
\item 增加一个 \textbf{cookies 字幕} case：由用户本机已登录浏览器 profile 提供登录态，期望 \texttt{selected\_source = official-cookies}
\item 继续保留当前 3 个 ASR-only Bilibili archetype，作为 mode routing 与 build 主链路 smoke set
\end{itemize}

\begin{riskbox}{为什么这是第一优先级}
因为系统设计从 README 到 wrapper runbook 都强调 subtitle-first。如果 live 层永远只跑 ASR fallback，那么当前最重要的设计约束实际上仍然停留在 mock 证明阶段。
\end{riskbox}

\subsection{建议 2：为 \texttt{transcripts.py} 建立 direct parser 回归}

建议新增 \texttt{tests/test\_transcripts.py}，至少覆盖以下场景：

\begin{itemize}
\item \texttt{parse\_srt\_text}：标准 SRT、多行字幕、空白块、逆序 end time
\item \texttt{parse\_vtt\_text}：带 \texttt{WEBVTT} 头、块间空行、多行正文
\item \texttt{parse\_json3\_payload}：YouTube JSON3 events 中空文本、零时长、多个 seg 拼接
\item \texttt{write\_transcript\_outputs} round-trip：写出的 \texttt{transcript.json / .srt / .txt} 是否与 payload 一致
\end{itemize}

这是高优先级，不是因为代码多，而是因为它处在“所有上游来源最终统一汇聚的单点”。一旦这里出错，后续写作、选图、时间窗对齐都会被污染。

\subsection{建议 3：为 adapter 层建立 contract tests}

建议新增 \texttt{tests/test\_adapters.py}，重点不是测网络，而是测 \textbf{contract}：

\begin{itemize}
\item \texttt{choose\_preferred\_track} 是否始终 manual 优先于 auto、中文优先于英文
\item \texttt{normalize\_metadata} 的缺省字段与 fallback 行为
\item \texttt{summarize\_best\_format} 的格式选择逻辑
\item \texttt{\_resolve\_subtitle\_sidecar} 是否正确定位 \texttt{.json3 / .srt / .vtt}
\item YouTube adapter 对 \texttt{YouTubeTranscriptApi} 返回对象的转换
\item Bilibili adapter 对 yt-dlp sidecar 的下载与解析约定
\end{itemize}

当前命令层测试把 adapter 大量 fake 掉，这对编排层是对的，但也让平台边界层成了一个明显空洞。

\subsection{建议 4：把 \texttt{probes/subtitles.py} 单独拉出来测}

建议新增 \texttt{tests/test\_subtitle\_probe.py}。至少应覆盖：

\begin{itemize}
\item 匿名成功时不再尝试 cookies
\item 匿名返回 \texttt{selected\_track} 但 payload 不可用时，仍继续尝试 cookies
\item cookies 也失败时，\texttt{selected\_track.source = none} 且 reason 合并清晰
\item selected payload 会被正确写入 \texttt{selected\_subtitle.json}
\end{itemize}

这类测试看起来和 \texttt{test\_commands.py} 有重叠，但价值在于：一旦 fallback 规则变复杂，probe 模块的 direct tests 会比命令级别 tests 更容易定位问题。

\subsection{建议 5：补齐 CLI 语义测试}

建议新增 \texttt{tests/test\_cli.py}，至少包括：

\begin{itemize}
\item 子命令 parser 是否接受预期参数
\item \texttt{--json} 输出是否是稳定 JSON
\item \texttt{CommandError} / \texttt{FileNotFoundError} 是否返回预期退出码
\item \texttt{build --note-tex}、\texttt{overview --sample-count} 等选项是否正确转发
\end{itemize}

当前系统的“真正用户入口”是 CLI，而不是内部 Python API。如果 CLI 不受保护，用户感知到的稳定性会低于内部代码质量。

\subsection{建议 6：补一个 wrapper/runtime contract sync test}

建议在 \path{/home/fanghaotian/Desktop/src/agent-basic-skill/tests/} 里新增专门针对 \texttt{video-note-render-pdf} 的 contract sync 测试，做的事情不是运行视频，而是比较：

\begin{itemize}
\item wrapper 文档中声明的必要 artifact
\item runtime \texttt{contracts.py} 里定义的 artifact 文件与目录
\item runtime 命令真实写回的关键字段，例如 \texttt{preflight.json.recommended\_mode}
\end{itemize}

原因很简单：当前 wrapper 测试已经证明“这个 skill 能装、能解析路径、文件都在”，但还没有证明“它说的话和 runtime 真的一致”。

\subsection{建议 7：把 build 的失败路径也锁住}

当前 \texttt{tests/test\_build.py} 已经覆盖 happy path 与外部 \texttt{note.tex}。建议继续增加：

\begin{itemize}
\item latexmk 返回非零时，\texttt{build\_summary.json} 是否如实标记失败
\item \texttt{note.tex} 缺失时，错误信息是否足够明确
\item 系统存在 \texttt{pdftoppm} 时，preview 图片是否正确落盘
\item 外部 \texttt{note.tex} 存在但没有生成 PDF 时，是否给出可定位的失败信号
\end{itemize}

\subsection{建议 8：没有 CI，就把“该跑什么”写死到 CI}

如果只谈新增测试而不谈执行入口，收益会被打折。建议的最小 CI 方案是：

\begin{itemize}
\item runtime repo：默认 PR 上跑 \texttt{uv run pytest}
\item wrapper repo：默认 PR 上跑 \texttt{uvx pytest tests/test\_wrapper\_resolution.py tests/test\_install\_skill.py tests/test\_skill\_repo.py}
\item live E2E：先不做强制门禁，但保留一个手动触发 workflow，至少跑公开字幕 canary 与当前 3 个 ASR smoke case
\end{itemize}

\section{结论与后续问题}

\subsection{已经有强证据支持的结论}

\begin{itemize}
\item 当前不是“缺少测试的原型仓库”，而是已有较清晰分层的测试体系
\item \texttt{video-note-pipeline} 当前工作副本的 runtime pytest 实测为 43/43 通过
\item \texttt{agent-basic-skill} 中与 \texttt{video-note-render-pdf} 直接相关的 wrapper pytest 实测为 17/17 通过
\item live E2E 已经覆盖 \texttt{visual-light}、\texttt{static-outline}、\texttt{board-heavy} 三类 archetype，并且 3/3 通过
\end{itemize}

\subsection{比较可能成立，但还没有被完全证明的结论}

\begin{itemize}
\item subtitle-first 的匿名字幕路径与 cookies 路径在真实 live 环境中长期稳定
\item adapter 面对 yt-dlp / YouTubeTranscriptApi 行为漂移时仍能保持 contract 稳定
\item wrapper 文档与 runtime contract 在未来演进中不会重新漂移
\end{itemize}

\subsection{仍然未知，值得下一轮验证的问题}

\begin{itemize}
\item 真实登录态 cookies case 是否能稳定进入自动化，而不把用户本地浏览器环境弄得太脆
\item transcript parser 在真实混杂字幕格式下的边角行为是否已经足够稳
\item 在没有 CI 的情况下，当前通过状态能否被持续维持，而不是依赖人工记忆
\end{itemize}

如果只允许我给一个最简路线图，我会建议按下面顺序推进：

\begin{enumerate}
\item 先把公开字幕 canary 和 cookies canary 加进 live E2E
\item 再补 \texttt{tests/test\_transcripts.py} 与 \texttt{tests/test\_adapters.py}
\item 接着补 \texttt{tests/test\_cli.py} 与 wrapper/runtime contract sync test
\item 最后把 runtime pytest 与 wrapper pytest 接进 CI
\end{enumerate}

\appendix

\section{来源清单}

\small
\begin{longtable}{p{1.2cm}p{1.5cm}p{3.2cm}p{2.1cm}p{6.5cm}}
\toprule
ID & 模态 & 标题 / 标签 & 可信度 & 链接或路径 \\
\midrule
\endfirsthead
\toprule
ID & 模态 & 标题 / 标签 & 可信度 & 链接或路径 \\
\midrule
\endhead
S1 & Code & runtime 依赖与 pytest 配置 & primary & \path{/home/fanghaotian/Desktop/src/video-note-pipeline/pyproject.toml} \\
S2 & Doc & runtime README 与 subtitle-first 说明 & primary & \path{/home/fanghaotian/Desktop/src/video-note-pipeline/README.zh-CN.md} \\
S3 & Code & runtime 命令编排测试 & primary & \path{/home/fanghaotian/Desktop/src/video-note-pipeline/tests/test_commands.py} \\
S4 & Code & cookies 测试 & primary & \path{/home/fanghaotian/Desktop/src/video-note-pipeline/tests/test_cookies.py} \\
S5 & Code & frames / mode routing 测试 & primary & \path{/home/fanghaotian/Desktop/src/video-note-pipeline/tests/test_frames.py} \\
S6 & Code & build 测试 & primary & \path{/home/fanghaotian/Desktop/src/video-note-pipeline/tests/test_build.py} \\
S7 & Code & transcriber probe 测试 & primary & \path{/home/fanghaotian/Desktop/src/video-note-pipeline/tests/test_transcriber_probe.py} \\
S8 & Code & live E2E harness 测试 & primary & \path{/home/fanghaotian/Desktop/src/video-note-pipeline/tests/test_live_e2e.py} \\
S9 & Data & 最新 live E2E 汇总 & primary & \path{/home/fanghaotian/Desktop/src/video-note-pipeline/reports/live-e2e/runs/bilibili-core-base-model-rerun/summary.json} \\
S10 & Code & cookies helper 当前实现 & primary & \path{/home/fanghaotian/Desktop/src/video-note-pipeline/src/video_note_pipeline/cookies.py} \\
S11 & Code & wrapper 路径解析测试 & primary & \path{/home/fanghaotian/Desktop/src/agent-basic-skill/tests/test_wrapper_resolution.py} \\
S12 & Code & wrapper 安装测试 & primary & \path{/home/fanghaotian/Desktop/src/agent-basic-skill/tests/test_install_skill.py} \\
S13 & Code & wrapper 包装结构测试 & primary & \path{/home/fanghaotian/Desktop/src/agent-basic-skill/tests/test_skill_repo.py} \\
S14 & Doc & wrapper 主说明 & primary & \path{/home/fanghaotian/Desktop/src/agent-basic-skill/skills/video-note-render-pdf/SKILL.md} \\
S15 & Doc & wrapper adapter contract & primary & \path{/home/fanghaotian/Desktop/src/agent-basic-skill/skills/video-note-render-pdf/references/adapter-contract.md} \\
S16 & Data & runtime pytest 审计日志 & primary & \path{/home/fanghaotian/Desktop/src/agent-basic-skill/reports/2026-03-26-video-note-testing-audit/runtime-pytest.log} \\
S17 & Data & wrapper pytest 审计日志 & primary & \path{/home/fanghaotian/Desktop/src/agent-basic-skill/reports/2026-03-26-video-note-testing-audit/wrapper-pytest.log} \\
S18 & Data & 源码模块到测试文件映射 & secondary & \path{/home/fanghaotian/Desktop/src/agent-basic-skill/reports/2026-03-26-video-note-testing-audit/source-module-test-map.txt} \\
S19 & Data & runtime dirty worktree 状态 & secondary & \path{/home/fanghaotian/Desktop/src/agent-basic-skill/reports/2026-03-26-video-note-testing-audit/runtime-git-status.txt} \\
\bottomrule
\end{longtable}

\section{附件说明}

本报告目录下还保留了以下审计附件，便于后续追溯：

\begin{itemize}
\item \path{runtime-pytest.log}
\item \path{wrapper-pytest.log}
\item \path{runtime-test-counts.csv}
\item \path{wrapper-test-counts.csv}
\item \path{source-module-test-map.txt}
\item \path{runtime-git-status.txt}
\item \path{wrapper-git-status.txt}
\end{itemize}

\end{document}
